%% kapa: exam=1 round=37 year=2026 period=1 form=A subject=economics subject_ko=경제학원론 seq=30 number=70
%% source: Q-Net 감정평가사 기출문제 (cst00309) artlSeq=5258924 file=raw/1008/5258924/2026년 제37회 감정평가사 1차 1교시.pdf
%% conversion: kordoc 4.24.51 -> tools/gen1.py (auto); review=proofread-sonnet; flags=none
\begin{question}{70}
일시적으로 수입 원유가격이 급상승하였다고 하자. 이것이 장기균형 상태에 있던 원유수입국에 미치는 영향을 총수요-총공급모형으로 분석할 때, 이에 관한 설명으로 옳은 것은? (단, 총수요곡선은 우하향, 단기 총공급곡선은 우상향, 장기 총공급곡선은 수직이다.)
\begin{choices}
  \item 단기적으로 물가가 상승하고 국민소득은 불변이다.
  \item 단기적으로 물가가 하락하고 국민소득은 증가한다.
  \item 장기적으로 물가는 원유가격 상승 이전 수준으로 돌아가고 국민소득은 감소한다.
  \item 장기적으로 물가와 국민소득 모두 원유가격 상승 이전 수준으로 돌아간다.
  \item 장ㆍ단기 모두 물가와 국민소득은 상승한다.
\end{choices}
\end{question}
\answer{4}
